%==============================================================================
%  Doing an undergraduate thesis in ML / NLP -- a practical guide
%  Konstantinos Skianis, University of Ioannina
%  Compile:  latexmk -pdf undergrad-thesis-ml-nlp-guide.tex
%==============================================================================
\documentclass[11pt,a4paper]{article}

\usepackage[T1]{fontenc}
\usepackage[utf8]{inputenc}
\usepackage{lmodern}
\usepackage{microtype}
\usepackage[a4paper,margin=2.3cm,top=2.4cm,bottom=2.4cm]{geometry}
\usepackage{xcolor}
\usepackage{graphicx}
\usepackage{amsmath,amssymb}
\usepackage{array,booktabs,tabularx,longtable,colortbl}
\usepackage{enumitem}
\usepackage{framed}
\usepackage{multicol}
\usepackage{fancyhdr}
\usepackage{caption}
\usepackage{url}
\usepackage{hyperref}
\usepackage{bookmark}

%------------------------------------------------------------------ colours ---
\definecolor{brand}{HTML}{1F4E79}      % deep blue
\definecolor{brandlight}{HTML}{2E7BB8}
\definecolor{accent}{HTML}{B5451B}     % rust
\definecolor{ink}{HTML}{22262B}
\definecolor{mutegray}{HTML}{5A6472}
\definecolor{rowtint}{HTML}{EFF4F9}
\definecolor{boxtint}{HTML}{F6F8FA}

\hypersetup{
  colorlinks=true,
  linkcolor=brand,
  urlcolor=brandlight,
  citecolor=brand,
  pdftitle={Doing an Undergraduate Thesis in Machine Learning or NLP},
  pdfauthor={Konstantinos Skianis},
  pdfsubject={Undergraduate thesis supervision, ML/NLP research practice},
  bookmarksnumbered=true
}
\urlstyle{sf}
% allow long URLs to break at punctuation instead of running into the margin
\makeatletter
\g@addto@macro\UrlBreaks{\do\/\do\-\do\.\do\_\do\?\do\&\do\=\do\:}
\makeatother
\Urlmuskip=0mu plus 2mu\relax

%----------------------------------------------------------------- headings ---
\makeatletter
\renewcommand\section{\@startsection{section}{1}{\z@}%
  {-3.2ex \@plus -1ex \@minus -.2ex}{1.6ex \@plus.2ex}%
  {\normalfont\Large\bfseries\color{brand}}}
\renewcommand\subsection{\@startsection{subsection}{2}{\z@}%
  {-2.6ex\@plus -1ex \@minus -.2ex}{0.9ex \@plus .2ex}%
  {\normalfont\large\bfseries\color{brandlight}}}
\renewcommand\subsubsection{\@startsection{subsubsection}{3}{\z@}%
  {-2.0ex\@plus -1ex \@minus -.2ex}{0.6ex \@plus .2ex}%
  {\normalfont\normalsize\bfseries\color{ink}}}
\makeatother

%-------------------------------------------------------------------- boxes ---
\newenvironment{cbox}[1]
 {\par\addvspace{\medskipamount}%
  \setlength{\FrameSep}{7pt}%
  \def\FrameCommand{{\color{brandlight}\vrule width 2.2pt}\hspace{9pt}}%
  \MakeFramed{\advance\hsize-\width \FrameRestore}%
  \small\parindent=0pt\noindent{\bfseries\color{brand}#1}\par\nobreak\smallskip
  \normalcolor\ignorespaces}
 {\par\endMakeFramed\addvspace{\medskipamount}}

\newenvironment{wbox}[1]
 {\par\addvspace{\medskipamount}%
  \setlength{\FrameSep}{7pt}%
  \def\FrameCommand{{\color{accent}\vrule width 2.2pt}\hspace{9pt}}%
  \MakeFramed{\advance\hsize-\width \FrameRestore}%
  \small\parindent=0pt\noindent{\bfseries\color{accent}#1}\par\nobreak\smallskip
  \normalcolor\ignorespaces}
 {\par\endMakeFramed\addvspace{\medskipamount}}

%------------------------------------------------------------------- tables ---
\newcolumntype{L}[1]{>{\raggedright\arraybackslash}p{#1}}
\renewcommand{\arraystretch}{1.22}
\newcommand{\thead}[1]{\textbf{\color{brand}#1}}

%-------------------------------------------------------------------- lists ---
\setlist[itemize]{leftmargin=1.25em,topsep=3pt,itemsep=2.5pt,parsep=0pt}
\setlist[enumerate]{leftmargin=1.5em,topsep=3pt,itemsep=2.5pt,parsep=0pt}
\setlist[description]{leftmargin=1.25em,topsep=3pt,itemsep=3pt,style=unboxed,
                      font=\normalfont\bfseries\color{ink}}

%------------------------------------------------------------------ headers ---
\pagestyle{fancy}
\fancyhf{}
\renewcommand{\headrulewidth}{0.3pt}
\fancyhead[L]{\footnotesize\color{mutegray}Undergraduate thesis in ML / NLP \quad\textbullet\quad K. Skianis}
\fancyhead[R]{\footnotesize\color{mutegray}\thepage}

\newcommand{\tool}[1]{\textbf{#1}}
\newcommand{\code}[1]{\texttt{\small #1}}

\begin{document}

%==============================================================================
\begin{titlepage}
\centering
\vspace*{2.0cm}
{\color{brandlight}\rule{\textwidth}{2pt}}\\[0.8cm]
{\Huge\bfseries\color{brand} Doing an Undergraduate Thesis\\[0.25cm] in Machine Learning or NLP\par}
\vspace{0.7cm}
{\LARGE\color{ink} A Practical Guide\par}
\vspace{0.5cm}
{\large\color{mutegray} Finding a topic, doing the work, writing it up, defending it\par}
\vspace{0.6cm}
{\color{brandlight}\rule{\textwidth}{2pt}}\\[1.1cm]

{\large\bfseries\color{ink} Konstantinos Skianis\par}
\vspace{0.18cm}
{\normalsize\color{mutegray} Assistant Professor of Natural Language Processing\\[2pt]
Department of Computer Science and Engineering\\[2pt]
University of Ioannina\par}
\vspace{0.25cm}
{\footnotesize\color{brandlight}\url{https://y3nk0.github.io}\par}
\vspace{1.1cm}

\begin{minipage}{0.85\textwidth}
\small\color{ink}
Written for final-year students starting a thesis in machine learning or
natural language processing. It is deliberately concrete: how to find a topic
that is neither trivial nor impossible, how to run experiments you can defend,
how to write a thesis that reads well, and what to do with it afterwards.
Nobody expects you to beat the state of the art. They expect you to do
something correctly, understand it completely, and explain it clearly.
\end{minipage}

\vfill
{\color{mutegray} October 2026}\\[0.3cm]
{\footnotesize\color{mutegray} Your department's own regulations --- length,
format, deadlines, grading --- always take precedence over anything here.}
\end{titlepage}

%==============================================================================
\tableofcontents
\thispagestyle{fancy}
\clearpage

%==============================================================================
\section{What a thesis is, and what it is not}

An undergraduate thesis is a supervised piece of work, typically 4--9 months
alongside your other courses, that shows you can take a research-level problem,
work on it methodically, and communicate the result. It is the first time most
students own a project end to end.

\begin{cbox}{What examiners actually look for}
\begin{itemize}
  \item A \textbf{clear question}, stated in one sentence, that your work
        answers.
  \item \textbf{Correct methodology} --- sane baselines, honest evaluation, no
        leakage, results you can reproduce.
  \item \textbf{Understanding.} You must be able to explain every design
        decision and every line you submit. This is what the defence tests.
  \item \textbf{Clear writing} --- a reader from your department should
        understand what you did and why.
  \item \textbf{Honest limitations.} Saying what did not work, and why, is a
        sign of maturity, not of failure.
\end{itemize}
\end{cbox}

\begin{wbox}{What is \emph{not} required}
Beating the state of the art. Inventing a new architecture. Training a large
model from scratch. A thesis that carefully reproduces, analyses and extends
existing work --- and explains what it found --- is a good thesis. A thesis
that chases a record and ends with one unexplained number is not.
\end{wbox}

\subsection{Scope: the single most common failure}

Almost every thesis that goes badly goes badly for the same reason: the topic
was too large, so the student spent six months getting a pipeline to run and
had nothing to say about it. Use this test:

\begin{itemize}
  \item Can you name the \textbf{dataset} and the \textbf{metric} today?
  \item Can you run \textbf{one baseline} on your hardware (or on free cloud
        GPUs) within \textbf{two weeks}?
  \item If your main idea fails completely, is there still a thesis ---
        a comparison, an analysis, an error study?
\end{itemize}

Three yeses mean the scope is right. Any no means narrow it with your
supervisor before you write code.

%==============================================================================
\section{Finding a topic: three routes}
\label{sec:topic}

You do not need an original idea to start. You need a \emph{well-defined
problem} with data and a way to measure success. There are three reliable
routes to one; pick whichever fits how you like to work --- but check what is
already on offer before you start inventing.

\begin{cbox}{Step zero: check what is already on offer}
Most supervisors keep a list of open topics --- on their web page, on the
department's thesis board, or in a shared document --- and the quickest good
thesis usually starts there. The order that works:
\begin{enumerate}
  \item \textbf{Look for advertised topics first.} Read the supervisor's page
        and the department's list, and ask the secretariat or older students
        where topics are posted.
  \item \textbf{Pick one and do a small piece of research before the first
        meeting:} read the two or three papers the topic points to, find the
        dataset and the metric, skim or run one existing implementation, and
        write down three questions. Arriving with that turns the first meeting
        from an introduction into a plan.
  \item \textbf{Only if you are given freedom}, or none of the advertised
        topics fits, use the three routes below to find your own.
\end{enumerate}
\end{cbox}

\begin{center}
\small
\begin{tabularx}{\textwidth}{L{3.0cm} X L{3.6cm}}
\toprule
\thead{Route} & \thead{What you do} & \thead{Best if} \\
\midrule
\rowcolor{rowtint}
\textbf{A. Shared task or challenge} & Enter an organised competition: the
data, the metric and the deadline already exist & You want structure and a
fast start \\
\textbf{B. Application thesis} & Take a method that works and apply it
carefully to a new domain, language or setting & You like building systems \\
\rowcolor{rowtint}
\textbf{C. Analysis or theoretical thesis} & Reproduce, probe, compare,
measure: explain \emph{why} methods behave as they do & You like understanding
over building \\
\bottomrule
\end{tabularx}
\end{center}

\subsection{Route A --- shared tasks, challenges and SemEval}
\label{sec:sharedtasks}

This is the strongest default for an undergraduate thesis, and it is
under-used. A \textbf{shared task} is a community competition: organisers
publish a precisely defined problem, training and development data, an official
metric, and an evaluation window. Everyone is evaluated on the same hidden test
set, and participants write a short \textbf{system description paper}.

\begin{cbox}{Why this is such a good fit}
\begin{itemize}
  \item \textbf{The hardest part is done for you.} Problem definition, data
        collection, annotation and the evaluation protocol --- the parts that
        sink inexperienced projects --- are settled.
  \item \textbf{You cannot fool yourself.} A hidden test set and a public
        leaderboard give you an honest measure of where you stand.
  \item \textbf{You get real baselines.} Organisers publish one, and past
        editions give you a whole table of comparisons for free.
  \item \textbf{There is a deadline}, which is the best project-management tool
        ever invented.
  \item \textbf{It can end in a publication.} System description papers are
        peer-reviewed, archival, and genuinely achievable for a strong
        undergraduate --- often your first paper, with your supervisor.
  \item \textbf{It is a portfolio item.} ``4th of 31 teams'' on a CV is
        concrete in a way that ``did a thesis on sentiment analysis'' is not.
\end{itemize}
\end{cbox}

\subsubsection*{Where to find them}

\begin{description}
  \item[SemEval] (\url{https://semeval.github.io}) --- the International
        Workshop on Semantic Evaluation, run annually with a dozen or more
        tasks: semantic similarity, sentiment and emotion, irony, humour,
        hate speech, factuality and hallucination detection, clinical and
        legal NLP, multilingual tasks, and more. Tasks are announced well
        ahead of the evaluation window and the workshop is co-located with a
        \mbox{*ACL} conference.
  \item[Workshop shared tasks at *ACL venues] --- most workshops at ACL,
        EMNLP, NAACL, COLING and EACL run their own task. Check
        \url{https://aclanthology.org} and the workshop lists on each
        conference site. Examples of recurring families: \textbf{WMT}
        (machine translation and evaluation metrics), \textbf{CoNLL} shared
        tasks, \textbf{BabyLM} (sample-efficient language modelling),
        \textbf{BioNLP}, \textbf{WNUT} (noisy user-generated text),
        \textbf{CheckThat!}~and other fact-checking labs.
  \item[CLEF labs and TREC tracks] --- \url{https://www.clef-initiative.eu},
        \url{https://trec.nist.gov}: information retrieval, biomedical
        question answering (\textbf{BioASQ}), argument mining, authorship,
        misinformation. Longer cycles, very well organised.
  \item[Competition tracks at ML conferences] --- the NeurIPS competition
        track and similar programmes at other venues: usually larger in scope,
        worth a look if you have compute.
  \item[Platforms] --- \textbf{Kaggle} (\url{https://www.kaggle.com}),
        \textbf{CodaBench / CodaLab}, \textbf{EvalAI},
        \textbf{AIcrowd}, \textbf{DrivenData}. Kaggle is excellent for
        learning, but a Kaggle competition alone is usually not a thesis:
        you still need a research question and an analysis chapter around it.
  \item[Leaderboards and benchmark hubs] --- Hugging Face
        (\url{https://huggingface.co}) hosts datasets, models and evaluation
        leaderboards; many tasks keep a permanent leaderboard open long after
        the official window closes.
\end{description}

\begin{wbox}{Check the dates before you commit}
Shared tasks run on fixed cycles --- registration, training data release,
development phase, a short evaluation window on the hidden test set, then the
system paper deadline --- and the exact dates change every edition. Read the
task's own page, write the four dates into your calendar, and confirm with your
supervisor that they fit your thesis deadline. An evaluation window that opens
after your submission date is a thesis you cannot finish.
\end{wbox}

\subsubsection*{How to run a shared task as a thesis}

\begin{enumerate}
  \item \textbf{Pick two or three candidate tasks} and read last year's
        \emph{task overview paper} plus three or four system papers. Thirty
        minutes each. You will immediately see what a good system looks like
        and where the easy gains are.
  \item \textbf{Reproduce the organisers' baseline first.} Before anything
        clever. If you cannot match the baseline, nothing else you do is
        interpretable.
  \item \textbf{Build your own validation split} and trust it. The leaderboard
        has limited submissions; your development set does not.
  \item \textbf{Keep an experiment log from run one} --- configuration, seed,
        score, what you changed, what you expected. This log becomes your
        results chapter.
  \item \textbf{Submit something early}, even a weak system. Submission
        pipelines always break the first time; find out in week three, not on
        the last day.
  \item \textbf{Write the system paper} with your supervisor. Then the thesis
        is the long version: the paper plus the background, the failed
        attempts, the ablations and the analysis that did not fit in four
        pages.
\end{enumerate}

\begin{cbox}{If the task window has closed}
Use it anyway. Past-edition data, metrics and leaderboards are usually public,
so you can run the whole thing retrospectively and compare against the
published table of participants. You lose the live leaderboard; you keep every
other advantage --- and you can report where your system \emph{would} have
ranked.
\end{cbox}

\subsection{Route B --- an application thesis}

Take a method that is known to work and apply it, carefully, somewhere it has
not been applied. The contribution is the careful application, the evaluation
and the analysis --- not the method.

Good sources of a target:

\begin{itemize}
  \item \textbf{A language or dialect with few resources.} Most NLP results
        are reported on English. Greek, and most other languages, have far
        fewer evaluated baselines --- building and \emph{properly evaluating}
        one for a specific task is a real contribution.
  \item \textbf{A domain with its own vocabulary:} legal, medical, financial,
        educational, agricultural, archival text; code; sensor or log data.
  \item \textbf{A deployment constraint:} the model must run on a laptop, a
        phone or a Raspberry Pi; must answer in under 100\,ms; must work with
        50 labelled examples; must run without sending data to a cloud API.
  \item \textbf{A real user with a real problem} --- a lab in another
        department, a university service, a local company, a public dataset
        nobody has modelled. Scope it tightly and get the data question
        answered in week one.
\end{itemize}

\begin{cbox}{What makes an application thesis rigorous rather than a demo}
\begin{itemize}
  \item At least two baselines, one of them embarrassingly simple (majority
        class, TF-IDF + logistic regression, nearest neighbour, a rule). If
        the simple baseline wins, that \emph{is} your finding --- report it.
  \item A fixed split made \emph{before} you look at results, and a test set
        touched once.
  \item Several seeds, with mean and spread. One run proves nothing.
  \item An error analysis: categorise 50--100 mistakes by hand. This is the
        single most valuable chapter in most application theses, and the one
        examiners remember.
  \item A cost column: parameters, training time, inference latency, hardware.
        Accuracy per euro and per second is a real axis.
\end{itemize}
\end{cbox}

\subsection{Route C --- an analysis or theoretical thesis}

You do not have to build something new to contribute. Explaining \emph{why}
existing methods behave as they do is harder and often more valuable.

\begin{itemize}
  \item \textbf{Reproduction and stress-testing.} Take a well-known recent
        result and check it: other seeds, other datasets, other languages,
        other hyperparameters. Reproducibility studies are a recognised genre
        with their own venues and tracks.
  \item \textbf{Ablation and attribution.} A paper proposes five components
        and reports one number. Which component actually matters? Very often
        the answer is ``not the one in the title''.
  \item \textbf{Robustness and failure modes.} What happens under typos,
        dialect, code-switching, domain shift, adversarial rewrites, longer
        inputs, class imbalance?
  \item \textbf{Probing and interpretability.} What does the representation
        encode? Where does the behaviour live in the network? Does the
        explanation match the mechanism?
  \item \textbf{Efficiency.} Distillation, quantisation, pruning, retrieval
        instead of parameters --- measured honestly on the same hardware.
  \item \textbf{Evaluation itself.} Does the metric measure what everyone
        assumes? Build a small, careful human evaluation and compare. Metric
        critiques are consistently under-supplied and over-needed.
  \item \textbf{Genuine theory} --- a proof of convergence, a bound, a
        complexity result, an equivalence between two formulations --- if you
        enjoy mathematics and have a supervisor to work with closely. Keep the
        claim small and complete rather than grand and sketched.
\end{itemize}

\subsection{How to find a gap without reading a thousand papers}

\begin{cbox}{Six places where gaps are written down explicitly}
\begin{enumerate}
  \item \textbf{``Limitations'' sections.} Every ACL-venue paper must have one.
        Authors list, in their own words, what their method cannot do. That
        list is a topic menu.
  \item \textbf{``Future work''} in the conclusion of recent papers and
        surveys.
  \item \textbf{Survey taxonomies.} Find the most recent survey in the area
        and look for the boxes with one or two entries --- that is where the
        field is thin.
  \item \textbf{Workshop calls for papers.} They list the open problems the
        community cares about this year, in bullet form.
  \item \textbf{Task overview papers.} The closing section of a shared-task
        overview says exactly what nobody solved.
  \item \textbf{Leaderboards that have plateaued}, or where a simple baseline
        is suspiciously close to the top --- a sign the task, the metric or
        the data deserves a second look.
\end{enumerate}
\end{cbox}

Two more practical sources: \textbf{GitHub issues} on popular repositories
(``does not work for language X'', ``cannot reproduce Table 2''), and
\textbf{reproducibility reports} that failed --- a failed reproduction is a
thesis-sized question.

\subsection{The one-page proposal}

Before you write any code, write one page and send it to your supervisor. It
takes an hour and saves months.

\begin{center}
\small
\begin{tabularx}{0.95\textwidth}{L{3.4cm} X}
\toprule
\thead{Field} & \thead{What to write} \\
\midrule
\rowcolor{rowtint} Question & One sentence, answerable, with a verb \\
Why it matters & Two sentences, no hype \\
\rowcolor{rowtint} Data & Name, size, licence, where it lives, is it public \\
Metric & The official one, plus what you will also report \\
\rowcolor{rowtint} Baselines & The trivial one and the strong one \\
Method & What you will try, in three bullets \\
\rowcolor{rowtint} Compute & Your laptop / lab GPU / Colab --- and whether it suffices \\
Milestones & Four dates: baseline, first result, freeze, draft \\
\rowcolor{rowtint} Risks \& fallback & What could fail, and the thesis you still have if it does \\
\bottomrule
\end{tabularx}
\end{center}

%==============================================================================
\section{Your toolkit, set up in week one}
\label{sec:tools}

\subsection{References: Zotero}

You will read 40--100 papers. Without a system you will re-find the same paper
five times and mis-cite one in your thesis.

\begin{itemize}
  \item \textbf{Zotero} (\url{https://www.zotero.org}), free, plus the
        \textbf{browser connector}: one click on an arXiv, ACL Anthology or
        Scholar page saves the metadata \emph{and} the PDF.
  \item \textbf{Better BibTeX} (\url{https://retorque.re/zotero-better-bibtex/}):
        stable citation keys you \emph{pin} (\code{devlin2019bert}), and
        \textbf{auto-export} of your library into the \code{refs.bib} of your
        thesis repository. Your bibliography is then always current and you
        never hand-edit BibTeX.
  \item Take a one-line note on every save: \emph{what it does, why I care}.
        In month six that is the difference between a library and a junkyard.
  \item Cite the \textbf{published version}, not the preprint, when one exists.
        Get clean BibTeX from the \textbf{ACL Anthology}
        (\url{https://aclanthology.org}) or \textbf{DBLP}
        (\url{https://dblp.org}) --- Google Scholar's is often sloppy about
        venue names and pages.
  \item Zotero also checks your library against Retraction Watch: a red
        banner on an item means it was retracted. Find out why before citing.
  \item \textbf{Unpaywall} (\url{https://unpaywall.org/products/extension}),
        a browser extension that finds the legal free copy of a paywalled
        paper --- essential if your university's access is patchy.
\end{itemize}

\subsubsection*{Getting a reference into your \code{.bib}}

The easiest route needs no setup at all:

\begin{enumerate}
  \item In Google Scholar, open \emph{Settings} $\rightarrow$
        \emph{Bibliography manager} and tick \emph{Show links to import
        citations into BibTeX}. One time only.
  \item A \textbf{Cite} link now appears under every result. Click it, then
        \textbf{BibTeX}, and copy the entry.
  \item Paste it into your \code{refs.bib} --- or into Zotero with
        \emph{File $\rightarrow$ Import from clipboard}, which keeps your
        library and your bibliography in step. Scholar can also export in
        bulk: star papers into \emph{My library}, select them, then
        \emph{Export} $\rightarrow$ \emph{BibTeX}.
\end{enumerate}

\begin{wbox}{Check whatever Scholar gives you}
Scholar's entries are often sloppy: \code{arXiv preprint arXiv:1810.04805} for
a paper published at NAACL, missing pages and editors, venue names that differ
between entries, mangled accents. A bibliography full of ``arXiv preprint'' is
the clearest sign that nobody checked. Cleaner sources, each one click away:
\begin{itemize}
  \item \textbf{ACL Anthology} --- a BibTeX box on every paper page, and a
        full \code{.bib} for each volume.
  \item \textbf{DBLP} --- per-paper BibTeX and a whole-author export.
  \item \textbf{arXiv} --- ``Export BibTeX citation'' on the abstract page.
  \item \textbf{Semantic Scholar}, ACM, IEEE and Springer pages all have a
        cite or export button.
  \item \textbf{Any DOI, from the terminal:}
        \code{curl -LH "Accept: application/x-bibtex" https://doi.org/<DOI>}
        returns the publisher's own entry.
\end{itemize}
\end{wbox}

Past twenty or so references, stop copy-pasting: the Zotero connector plus
Better BibTeX auto-export does all of this for you and keeps your keys stable.

Two things every tool gets wrong unless you tell it: protect capitals in
titles (\code{\{BERT\}}, \code{\{Greek\}}) or the bibliography style will
lowercase them, and write authors as \code{Last, First and Last, First}.
Before you submit, run \code{bibtex-tidy}
(\url{https://github.com/FlamingTempura/bibtex-tidy}) over \code{refs.bib}
to format, sort and deduplicate it, and check that every DOI resolves ---
examiners now look for invented references.

\subsection{Writing: LaTeX and Overleaf}

Write the thesis in LaTeX unless your department requires otherwise. Learn it
in a weekend and it never costs you again. Typst is nice, but your
department's template and every venue's are LaTeX.

\begin{itemize}
  \item \textbf{Overleaf} (\url{https://www.overleaf.com}) --- browser-based,
        nothing to install, real-time co-editing and comments. Your supervisor
        can read and annotate the draft without exchanging files. Check whether
        your university has an institutional licence before paying. It now has
        a built-in, metered AI assistant as well; the rules in \S\ref{sec:ai}
        apply to it exactly.
  \item \textbf{Prism} (\url{https://prism.openai.com}) --- OpenAI's
        AI-native LaTeX workspace \textbf{shuts down on 1 January 2027}. Do
        not start your thesis there. If you already have projects, export them
        now: exports do not include comments, version history, collaborators
        or linked integrations.
  \item \textbf{Locally:} \code{git} + \code{latexmk -pdf -pvc} + VS Code.
        Write \textbf{one sentence per line} --- it makes diffs readable and
        revisions painless; \code{latexdiff} shows your supervisor exactly
        what changed between drafts.
  \item Start from your department's thesis template if one exists; otherwise
        a clean \code{report} or \code{book} class is fine. Ask early about
        required title-page wording, logos and declarations.
\end{itemize}

\subsection{Code, data and compute}

\begin{itemize}
  \item \textbf{One git repository} for the whole thesis: code, configs,
        notes, and the LaTeX. Commit daily. ``It was working yesterday'' is
        only recoverable if yesterday is in git.
  \item \textbf{Free GPUs:} Google Colab and Kaggle Notebooks both give usable
        GPU time; Kaggle sessions are longer and more predictable. Ask your
        supervisor about lab machines or university clusters --- and ask in
        week one, because accounts take time.
  \item \textbf{Hugging Face} (\url{https://huggingface.co}) for pretrained
        models and datasets. Check the \textbf{licence} of every model and
        dataset you use, and say in the thesis which you used and under what
        terms.
  \item \textbf{Experiment tracking:} Weights \& Biases, MLflow or
        TensorBoard. Log the config, the metrics and the git commit for every
        run. ``Which run produced Table 4?'' must have an answer in month six.
  \item \textbf{Data protection:} if your data involves people --- scraped
        social media, user records, recordings, surveys --- raise it with your
        supervisor before you collect anything. Licensing and GDPR are not
        optional, and ``I found it online'' is not a licence.
\end{itemize}

%==============================================================================
\section{Finding and judging related work}

\subsection{Google Scholar, properly}

\begin{center}
\small
\begin{tabularx}{\textwidth}{L{4.8cm} X}
\toprule
\thead{Query / feature} & \thead{What it does} \\
\midrule
\rowcolor{rowtint}
\code{"hate speech detection"} & Exact phrase. Always quote technical terms. \\
\code{greek OR hellenic} & \code{OR} must be capitalised. \\
\rowcolor{rowtint}
\code{-survey} & Excludes a term. \\
\code{author:"K Skianis"} & Author search. \\
\rowcolor{rowtint}
\code{source:"SemEval"} & Restricts to a publication whose name contains the
phrase --- the way to search inside a venue or a task. \\
\code{intitle:transformer} & Term must appear in the title: a brutal, effective
precision filter. \\
\rowcolor{rowtint}
\textbf{Custom range} + \textbf{Sort by date} & This month's work rather than
2019's most-cited. \\
\textbf{Cited by} $\rightarrow$ \textbf{search within citing articles} &
Search \emph{inside} the papers citing a key work. The fastest way to find who
applied a method to \emph{your} setting. \\
\rowcolor{rowtint}
\textbf{All $N$ versions} & Finds the free PDF and the published version. \\
\bottomrule
\end{tabularx}
\end{center}

Set up two or three \textbf{alerts} --- one for your topic, one for citations
of the paper your work builds on --- and triage them into Zotero once a week.

Beyond Scholar: the \textbf{ACL Anthology} (every *ACL paper, free, full-text
search), \textbf{Semantic Scholar} (TLDR summaries, citation intent),
\textbf{Connected Papers} (\url{https://www.connectedpapers.com}, a visual
neighbourhood from one seed paper --- excellent in week one),
\textbf{OpenReview} (\url{https://openreview.net}, where you can read the
actual \emph{reviews} of papers), and \textbf{paperswithcode.co}.

\textbf{Scholar Labs.} \url{https://scholar.google.com/scholar_labs/search} is a search mode
inside Scholar, open to all signed-in users: type the research question itself
rather than keywords (``does data augmentation help low-resource Greek
NER?'') and it returns papers that answer it, with a sentence on how each one
does and follow-up questions to narrow in. Good for the first week, when you
do not know the vocabulary yet; then read the papers --- the explanation is a
lead, not a citation.

\subsection{Reading efficiently: three passes}

\begin{cbox}{Keshav's method, which will save you weeks}
\textbf{Pass 1} (5 minutes): title, abstract, introduction, section headings,
conclusions --- then decide whether to continue. \textbf{Pass 2} (1 hour):
figures, tables and the method, to the level where you could summarise it.
\textbf{Pass 3} (several hours, rare): re-derive the work. Give pass 3 to the
five papers your thesis depends on, pass 1 to a hundred.
\end{cbox}

\subsection{Is this work worth building on?}

\begin{itemize}
  \item \textbf{Citations, normalised by age.} A 2026 paper with 40 citations
        may be hotter than a 2019 paper with 400. Scholar's per-year bar chart
        shows whether interest is accelerating or dead.
  \item \textbf{Who cites it} matters more than how many. Ten citations from
        strong groups beat 200 from unrelated surveys.
  \item \textbf{GitHub signals:} last commit date, number of contributors, who
        answers issues, whether releases and pinned dependencies exist, and
        whether checkpoints are provided. \emph{Search the issue tracker for
        ``reproduce''} before you trust a repository's numbers --- do this
        every time.
  \item \textbf{Stars measure marketing, not correctness.} The only real test
        is running it yourself on a task you understand. Use
        \url{https://star-history.com} to tell a one-day launch spike from
        steady adoption.
  \item \textbf{Licence.} Research-only or non-commercial terms can restrict
        what you are allowed to do and publish. Check before you invest a
        month.
\end{itemize}

%==============================================================================
\section{Using an AI assistant honestly}
\label{sec:ai}

An AI assistant will speed up the mechanical parts of your thesis
considerably. It can also end your degree if you use it to replace the thinking
the thesis is meant to demonstrate. Both of those are true at the same time.

\subsection{Set it up}

\begin{itemize}
  \item A coding assistant in your thesis repository, with a short project
        memory file stating your conventions: Python version, how to launch
        training, where data lives, plotting style.
  \item Connect it to your filesystem and, if you can, to your Zotero library
        and literature APIs --- an assistant that reads your actual PDFs is a
        different tool from one that guesses.
  \item Keep a dated \code{notes/} folder it can read and write: decisions,
        failed runs, hyperparameters.
\end{itemize}

\subsection{Use it for}

\begin{multicols}{2}
\begin{itemize}
  \item Explaining unfamiliar notation and walking through a derivation.
  \item Getting someone else's repository to run.
  \item Boilerplate: data loaders, plotting, evaluation scripts, Slurm or
        Colab setup.
  \item Refactoring code before it becomes unmaintainable.
  \item LaTeX debugging and table formatting.
  \item Tightening prose you have already written.
  \item Acting as a hostile examiner on your draft --- and, for a free first
        review of the PDF, Google's Deep Review (box below).
\end{itemize}
\end{multicols}

\begin{cbox}{A free first review of your draft: Deep Review}
Google's \textbf{Deep Review} (\url{https://research.google/deepreview/}) reads a PDF
like a careful examiner: it checks each section for mathematical errors,
logical gaps and unsupported claims, then the whole for consistency. Two
reviews a week are free, and Google states that uploads are not used to train
its models. Run it on your complete draft before your supervisor reads it,
then read every flag: fix what is wrong and note why you reject the rest. It
is one more hostile reader, not a verdict --- and never upload someone else's
unpublished work.
\end{cbox}

\subsection{The rules}

\begin{wbox}{Five rules, and the first one is not negotiable}
\begin{enumerate}
  \item \textbf{You must be able to explain every line you submit.} At the
        defence you will be asked why a function is there and what a number
        means. ``The assistant wrote it'' ends the conversation badly. If you
        do not understand it, do not submit it.
  \item \textbf{Never cite a paper you have not opened.} Models invent
        plausible references --- right-looking authors, venue and year, all
        fabricated. Every entry in your bibliography must come from Zotero,
        DBLP or the Anthology after you have read at least the abstract.
        Fabricated references are the fastest way to fail a thesis.
  \item \textbf{Never let it write your results or your analysis.} Numbers
        come from runs you did; interpretation comes from you.
  \item \textbf{Follow your university's policy} on AI assistance and declare
        your use if a declaration is required. Ask your supervisor what the
        rule is in your department --- do not guess.
  \item \textbf{Do not paste confidential or personal data} --- someone's
        unpublished manuscript, data under an NDA, or personal records ---
        into a cloud tool.
\end{enumerate}
\end{wbox}

%==============================================================================
\section{Doing the work}

\subsection{The order that works}

\begin{enumerate}
  \item \textbf{Get the data and look at it by hand.} Open fifty examples.
        Count the labels. Find the duplicates and the empty rows. One hour
        here saves weeks of confusion later.
  \item \textbf{Build the dumbest possible baseline} and measure it. Majority
        class, TF-IDF + logistic regression, a rule. This number anchors
        everything that follows.
  \item \textbf{Reproduce the strong baseline} from the paper or the task.
  \item \textbf{Only then} try your idea --- one change at a time, each
        measured.
  \item \textbf{Ablate.} Remove each piece of your method and show what it was
        worth. Ablations always take longer than the main result; start them
        early.
  \item \textbf{Analyse the errors} by hand, and categorise them.
\end{enumerate}

\subsection{Mistakes that cost marks}

\begin{wbox}{The usual suspects}
\begin{itemize}
  \item \textbf{Test-set leakage} --- tuning on test, duplicates shared between
        train and test, preprocessing or feature selection fitted on the full
        dataset before splitting.
  \item \textbf{One seed.} A single run with a $+0.4$ improvement is noise.
        Report mean $\pm$ standard deviation over at least three seeds.
  \item \textbf{Unfair comparison} --- your method tuned for a week, the
        baseline run with defaults for an hour. Examiners spot this
        immediately.
  \item \textbf{Accuracy on imbalanced data.} Report macro-F1, per-class
        results, or a confusion matrix.
  \item \textbf{Results that cannot be regenerated} because the configuration
        lived in your shell history.
  \item \textbf{Plots without axis labels, units or a caption} that says what
        the reader should conclude.
\end{itemize}
\end{wbox}

\subsection{Keep a results table from day one}

One spreadsheet or CSV: date, git commit, config file, seed, dev score, test
score (only when you finally run it), and a one-line note on what changed and
why. Your results chapter is then a formatting exercise rather than an
archaeology project, and the final table regenerates from logged numbers with a
script.

\begin{cbox}{Negative results are results}
``I tried X; it did not help; here is the evidence and here is why I think it
failed'' is a legitimate and valuable thesis finding --- provided the
experiment was fair and the analysis is honest. Many strong theses are built
exactly this way. Keep the failed experiments in your log; do not delete them.
\end{cbox}

%==============================================================================
\section{Writing the thesis}

Start writing in month one, not month six. Writing is how you discover that
your story has a hole, and holes are cheap to fix early.

\subsection{A structure that works}

\begin{center}
\small
\begin{tabularx}{\textwidth}{L{3.5cm} X L{2.0cm}}
\toprule
\thead{Chapter} & \thead{What goes in it} & \thead{Rough size} \\
\midrule
\rowcolor{rowtint}
Abstract & Problem, what you did, headline result, in one paragraph & 1 page \\
Introduction & Context, the question, your contributions as a bullet list,
thesis outline & 3--5 \\
\rowcolor{rowtint}
Background & Only the concepts the reader needs for your chapters & 5--10 \\
Related work & Organised by \emph{theme}, not one paragraph per paper; end by
saying where your work sits & 5--8 \\
\rowcolor{rowtint}
Method / system & What you built, precisely enough to reimplement & 5--10 \\
Experiments & Data, splits, metrics, baselines, setup, then results tables &
8--15 \\
\rowcolor{rowtint}
Analysis & Ablations, error analysis, what surprised you & 4--8 \\
Conclusion & What you found, limitations, future work & 2--3 \\
\rowcolor{rowtint}
References, appendix & Bibliography; extra tables, prompts, hyperparameters &
--- \\
\bottomrule
\end{tabularx}
\end{center}

\emph{Follow your department's required length and format; the sizes above are
a sane default, not a rule.}

\subsection{Writing advice that actually helps}

\begin{itemize}
  \item \textbf{Write the figures and tables first}, then the text that
        explains them. If a table needs three paragraphs of excuses, the
        experiment is wrong, not the prose.
  \item \textbf{Every chapter opens with one sentence saying what it does} and
        closes with one saying what it showed.
  \item \textbf{Say what you did, not what one might do.} Past tense, active
        voice, short sentences. ``We tokenise with X'' beats ``it was deemed
        appropriate to employ tokenisation''.
  \item \textbf{Define every symbol} the first time it appears, and keep
        notation consistent across chapters.
  \item \textbf{Captions are self-contained.} A reader who looks only at your
        figures should still understand the thesis.
  \item \textbf{Cite properly and quote almost never.} Copying sentences ---
        from papers, from websites, or from a model --- without quotation
        marks and a citation is plagiarism, and it is checked. Paraphrase in
        your own words and cite the source.
  \item \textbf{Use the right citation command.}
        \code{\textbackslash citet\{\}} when the authors are part of the
        sentence (``Devlin et al.\ (2019) show''),
        \code{\textbackslash citep\{\}} when the citation is an aside
        (``(Devlin et al., 2019)''); with biblatex,
        \code{\textbackslash textcite} and \code{\textbackslash parencite}.
        Mixing them is the most common formatting slip in student theses.
  \item \textbf{Give your supervisor a complete draft}, not a chapter, at
        least three weeks before the deadline. Feedback on an incomplete draft
        is feedback on the wrong problem.
\end{itemize}

%==============================================================================
\section{The defence}

Most departments ask for a 15--25 minute presentation plus questions. It is not
an interrogation; it is you explaining your work to people who want you to
succeed.

\begin{itemize}
  \item \textbf{Structure:} problem (2 slides) $\rightarrow$ why it is hard
        (1) $\rightarrow$ what you did (3--5) $\rightarrow$ results (3--4)
        $\rightarrow$ analysis and limitations (2) $\rightarrow$ conclusion
        (1). One idea per slide.
  \item \textbf{Rehearse out loud, timed, twice.} Once alone, once in front of
        a person. Speaking it aloud exposes everything you do not actually
        understand.
  \item \textbf{A short demo} --- 60 seconds, recorded as a fallback --- is
        worth a hundred words if you built a system. Never rely on live
        internet.
  \item \textbf{Know your numbers:} dataset sizes, your best score, the
        baseline's score, and how long training took.
  \item \textbf{Expect these questions:} Why this baseline? What happens if
        the data changes? Why does it fail on class X? What would you do with
        six more months? What did you find hardest? Is the improvement
        significant?
  \item \textbf{``I do not know'' is an acceptable answer} when followed by
        ``but here is how I would find out''. Guessing is not.
\end{itemize}

%==============================================================================
\section{A timetable}

A six-month plan; stretch or compress it to your own deadline, and put every
milestone in a calendar.

\begin{center}
\small
\begin{tabularx}{\textwidth}{L{2.2cm} X}
\toprule
\thead{Month} & \thead{What is done by the end of it} \\
\midrule
\rowcolor{rowtint}
\textbf{1} & Topic agreed and one-page proposal signed off; tools installed
(Zotero, git, Overleaf, assistant); data downloaded and inspected; 10--15
papers read; trivial baseline running. \\
\textbf{2} & Strong baseline reproduced; experiment logging in place;
background and related-work chapters drafted; shared-task dates in the
calendar (if Route A). \\
\rowcolor{rowtint}
\textbf{3} & Your method implemented; first real results; method chapter
drafted. \\
\textbf{4} & Ablations and seeds; results tables generated by script; first
full draft of the experiments chapter. \\
\rowcolor{rowtint}
\textbf{5} & Error analysis; analysis chapter; \textbf{complete draft} to your
supervisor; results frozen. \\
\textbf{6} & Revisions, polish, references checked one by one, defence slides,
two rehearsals, submission. \\
\bottomrule
\end{tabularx}
\end{center}

\begin{cbox}{Weekly rhythm with your supervisor}
Meet weekly or fortnightly, and send five bullets the evening before:
\textbf{what I did} $\cdot$ \textbf{what I found} $\cdot$ \textbf{what is
blocked} $\cdot$ \textbf{what I propose next} $\cdot$ \textbf{what I need from
you}. Keep a shared running document. Supervisors help far more effectively
when the question is specific, and ``I am stuck'' sent on Tuesday is worth
three weeks of silence.
\end{cbox}

%==============================================================================
\section{After the thesis}

\begin{itemize}
  \item \textbf{Publish it if it is good.} A shared-task system paper, a
        workshop paper, or a student research workshop submission (most *ACL
        conferences have one) is realistic for a strong thesis --- ask your
        supervisor. A paper transforms a thesis from coursework into a
        credential.
  \item \textbf{Clean the repository} and make it public: a README that
        explains how to reproduce your main table, the environment file, and
        a licence. This is the single best item on an entry-level CV.
  \item \textbf{Put up a page} --- \code{username.github.io} with the
        AcademicPages template takes an afternoon: a short bio, your thesis
        PDF, your code, your CV. Employers and MSc admissions committees
        search your name.
  \item \textbf{Ask for a reference letter while everything is fresh}, and
        give your supervisor your thesis abstract, your results and your plans
        to write from.
  \item \textbf{If you are considering a PhD}, the companion guide
        \emph{Starting a PhD in Machine Learning or NLP} covers venues,
        deadlines and the research habits that follow from here.
\end{itemize}

%==============================================================================
\section{Checklists}

\subsection*{Week 1}
\begin{multicols}{2}
\begin{itemize}
  \item Supervisor meeting slot fixed
  \item One-page proposal written and agreed
  \item Zotero + connector + Better BibTeX installed
  \item Git repository created, with \code{notes/} and a README
  \item Overleaf (or local LaTeX) project started
  \item Data downloaded, licence checked, 50 examples read by hand
  \item Compute sorted: lab machine, cluster account, or Colab/Kaggle
  \item Deadlines in the calendar, including any task dates
\end{itemize}
\end{multicols}

\subsection*{Before you freeze results}
\begin{itemize}
  \item Trivial baseline, strong baseline and your method, all on the same
        splits
  \item At least three seeds, mean $\pm$ std reported
  \item Test set used once; no tuning on it
  \item Every table and figure regenerated by a script from logged results
  \item Ablations done; error analysis categorised
  \item Code runs from a clean clone, following your own README
\end{itemize}

\subsection*{Before you submit}
\begin{itemize}
  \item Every citation checked against the real paper and every DOI
        resolves --- no invented references; \code{bibtex-tidy} run
  \item Department's format, length and declaration requirements met
  \item Figures legible in print and greyscale; captions self-contained
  \item Limitations section written honestly
  \item A Deep Review pass on the final PDF; every flag fixed or consciously
        rejected
  \item Spell-check and a full read-aloud pass
  \item A complete draft seen by your supervisor at least three weeks earlier
  \item Slides rehearsed twice, timed, with a recorded demo as a fallback
\end{itemize}

%==============================================================================
\appendix
\section{Link directory}

\begin{longtable}{L{4.3cm} L{10.5cm}}
\toprule
\thead{What} & \thead{Where} \\
\midrule
\endfirsthead
\toprule
\thead{What} & \thead{Where} \\
\midrule
\endhead
\multicolumn{2}{l}{\textbf{\color{brand}Shared tasks and challenges}}\\
SemEval & \url{https://semeval.github.io} \\
ACL Anthology (workshops) & \url{https://aclanthology.org} \\
CLEF initiative & \url{https://www.clef-initiative.eu} \\
TREC & \url{https://trec.nist.gov} \\
Kaggle & \url{https://www.kaggle.com} \\
CodaBench / CodaLab & \url{https://www.codabench.org} \\
EvalAI & \url{https://eval.ai} \\
AIcrowd & \url{https://www.aicrowd.com} \\
DrivenData & \url{https://www.drivendata.org} \\
Conference deadlines & \url{https://paperswithcode.co/ai-deadlines} \\
 & \url{http://aideadlines.org/?sub=ML,CV,NLP} \\
\midrule
\multicolumn{2}{l}{\textbf{\color{brand}Tools}}\\
Zotero & \url{https://www.zotero.org} \\
Better BibTeX & \url{https://retorque.re/zotero-better-bibtex/} \\
bibtex-tidy & \url{https://github.com/FlamingTempura/bibtex-tidy} \\
Overleaf & \url{https://www.overleaf.com} \\
Prism (closes 1 Jan 2027) & \url{https://prism.openai.com} \\
Unpaywall & \url{https://unpaywall.org/products/extension} \\
Deep Review (Google) & \url{https://research.google/deepreview/} \\
Hugging Face & \url{https://huggingface.co} \\
Google Colab & \url{https://colab.research.google.com} \\
Kaggle Notebooks & \url{https://www.kaggle.com/code} \\
Weights \& Biases & \url{https://wandb.ai} \\
\midrule
\multicolumn{2}{l}{\textbf{\color{brand}Literature}}\\
Google Scholar & \url{https://scholar.google.com} \\
Scholar Labs & \url{https://scholar.google.com/scholar_labs/search} \\
Semantic Scholar & \url{https://www.semanticscholar.org} \\
DBLP & \url{https://dblp.org} \\
OpenReview & \url{https://openreview.net} \\
Connected Papers & \url{https://www.connectedpapers.com} \\
Papers with Code & \url{https://paperswithcode.co} \\
\midrule
\multicolumn{2}{l}{\textbf{\color{brand}Afterwards}}\\
AcademicPages template & \url{https://github.com/academicpages/academicpages.github.io} \\
The PhD companion guide & \url{https://y3nk0.github.io/guides/} \\
\bottomrule
\end{longtable}

\section{One-page proposal template}

\begin{quote}\small
\textbf{Title.} Working title, one line.\\[3pt]
\textbf{Question.} The one sentence this thesis answers.\\[3pt]
\textbf{Motivation.} Two sentences: who cares, and why it is not already
solved.\\[3pt]
\textbf{Route.} Shared task / application / analysis (\S\ref{sec:topic}).\\[3pt]
\textbf{Data.} Name, size, source, licence, splits.\\[3pt]
\textbf{Metric.} Official metric, plus anything else you will report.\\[3pt]
\textbf{Baselines.} The trivial one; the strong one; where their numbers come
from.\\[3pt]
\textbf{Method.} Three bullets on what you will try.\\[3pt]
\textbf{Compute.} Hardware, and the largest model it allows.\\[3pt]
\textbf{Milestones.} Baseline by \rule{1.6cm}{0.4pt}; first result by
\rule{1.6cm}{0.4pt}; freeze by \rule{1.6cm}{0.4pt}; full draft by
\rule{1.6cm}{0.4pt}.\\[3pt]
\textbf{Risks.} What could fail.\\[3pt]
\textbf{Fallback.} The thesis you still have if it does.
\end{quote}

\vspace{0.5cm}
\begin{cbox}{The shortest version of all of this}
Pick a problem that already has data and a metric --- a shared task is the
easiest way to get one. Run the dumb baseline first. Change one thing at a
time and write down what happened. Understand everything you submit. Write
from month one, show your supervisor a complete draft early, and be honest
about what did not work. That is a good thesis.
\end{cbox}

\end{document}
